\chapter{Installation}\label{ch:install}

\section{Requirements}

\begin{itemize}
\item A C compiler with OpenMP, MPI, LAPACK and GSL.
\item PETSc and SLEPc, built with the same MPI. They are needed by phases
  2a and 2b. 64-bit indices are needed only for very large polyads
  (section~\ref{sec:intsize}).
\item ScaLAPACK for the distributed diagonalisation, and ELPA, which is needed to build phase 2b and is the fastest eigensolver in the tests (\opt{eigensolver}~=~3, Chapter~\ref{ch:tut-pyrene}). ELPA must link the same 64-bit BLAS, LAPACK and ScaLAPACK as the code: check with \texttt{ldd libelpa.so} that no 32-bit library (such as \texttt{libscalapack-openmpi} or \texttt{libopenblas} without the 64) comes in.
\item PARI/GP, for the optional search of hidden resonances
  (\opt{diophantineiter}).
\item The \texttt{oct\_parser} library, which reads the input file. It is
  included as the submodule \file{external/oct-parser}.
\end{itemize}

\section{Building}

\begin{lstlisting}
git clone --recursive <repository URL> AnharmoniCaOs
cd AnharmoniCaOs
./configure --with-petsc=PREFIX --with-slepc=PREFIX ...
make
\end{lstlisting}
\opt{./configure --help} lists the options that point to each library.
Building produces three programs:
\begin{itemize}
\item \prog{thermal\_anharmonicity\_fase1};
\item \prog{thermal\_anharmonicity\_fase2a};
\item \prog{thermal\_anharmonicity\_fase2b}.
\end{itemize}

\section{32-bit and 64-bit indices}\label{sec:intsize}

The integers passed to BLAS, LAPACK, ScaLAPACK, BLACS and ELPA have the
type \texttt{PetscBLASInt}. PETSc defines it as 32 or 64 bits, consistently
with the BLAS and LAPACK it was built with (\texttt{PETSC\_USE\_64BIT\_BLAS\_INDICES}
in \file{petscconf.h}). It is a different type from \texttt{PetscInt},
which is set by the PETSc option \texttt{--with-64-bit-indices}.

\begin{itemize}
\item \emph{When 64 bits are needed.} With 32-bit indices, a matrix or a
  distributed array passed to these libraries must have fewer than
  $2^{31}\approx 2.1\times10^{9}$ elements. A dense polyad of $N$ states
  has $N^2$ of them, so a matrix held whole reaches the limit at
  $N \approx 46\,000$. In distributed runs each process holds only part of
  the matrix, but the workspace sizes can also exceed the limit. Larger
  polyads need PETSc and all these libraries built with 64-bit
  integers (ILP64). Below this size the code builds and works with the
  usual 32-bit (LP64) libraries.
\item \emph{Consistency.} All of BLAS, LAPACK, ScaLAPACK, BLACS and ELPA
  must use the same integer size as \texttt{PetscBLASInt}. A library of the
  other size links without errors, but receives indices of the wrong width:
  the results are wrong, or the program crashes.
\item \emph{What \prog{configure} does.} \prog{configure}
  reads the size of \texttt{PetscBLASInt}. For BLAS, LAPACK and ScaLAPACK
  not given with \opt{--with-\ldots} or \texttt{NAME\_LIBS}, it tries the
  pkg-config modules available (for example \texttt{blas64},
  \texttt{lapack64}, \texttt{scalapack64-openmpi} on Debian and Ubuntu,
  or \texttt{mkl-dynamic-ilp64-gomp} for MKL, tried first when
  \texttt{MKLROOT} is set) and keeps the first one that passes a small
  test of its integer size; it lists the ones it skipped. It then tests
  the libraries actually used, including those given by the user, and
  stops if one of them uses the other size, and names it.
  \opt{--disable-int-check} turns the error into a warning. BLACS and ELPA
  are not checked: make sure they were built for the same integer size
  (for ELPA, a 64-bit build is usually configured with
  \texttt{--enable-64bit-integer-math-support}).
\end{itemize}

\paragraph{Licences of the libraries.} \ac{} is distributed under the GNU
GPL version~3. The libraries it links are free software compatible with it:
\begin{itemize}
\item GSL is under the GPL;
\item PETSc, SLEPc, ScaLAPACK and OpenBLAS are under BSD-style licences;
\item ELPA is under the LGPL;
\item the parser is from Octopus, under the GPL.
\end{itemize}
Intel MKL (oneMKL) is the exception: it is proprietary, distributed under
the Intel Simplified Software License, which allows its use and
redistribution without fees~\cite{intel-mkl-faq}, but it is not
GPL-compatible software. Building \ac{} with MKL for your own use is not a
problem. Whoever distributes executables linked with MKL must make sure
that they may do so. Free libraries with similar performance are
available: OpenBLAS, and the 64-bit ScaLAPACK of the Linux distributions
(\texttt{libscalapack64}).

\section{Using ELPA}\label{sec:elpa}

With \opt{eigensolver}~=~3 phase 2b diagonalises the polyads with ELPA, the
fastest of the eigensolvers in the tests of Chapter~\ref{ch:tut-pyrene}.
The code sets ELPA's solver with \opt{elpasolver} (1 one-stage, the
default; 2 two-stage) and the block size of the distributed matrices with
\opt{elpablock} (64).

\paragraph{What the installed ELPA contains.} \file{elpa2\_print\_kernels},
in the \file{bin} directory of the ELPA installation, lists the kernels of
the two-stage solver that were compiled in, and the default one. The
header \file{elpa/elpa\_configured\_options.h}, in the \file{include}
directory, says whether the GPU versions were built
(\texttt{ELPA\_WITH\_NVIDIA\_GPU\_VERSION 1}). \texttt{ldd libelpa.so}
shows the BLAS, LAPACK and ScaLAPACK it uses: they must be the 64-bit ones
of the code (Section~\ref{sec:intsize}).

\paragraph{Choosing at run time.} ELPA reads environment variables
\texttt{ELPA\_FORCE\_}\emph{option}, which override what the code sets,
\opt{elpasolver} included. With MPI they must be passed to every process,
for example
\begin{verbatim}
mpirun -np 1 -x ELPA_FORCE_solver=ELPA_SOLVER_2STAGE \
       -x ELPA_FORCE_real_kernel=ELPA_2STAGE_REAL_AVX2_BLOCK4 \
       thermal_anharmonicity_fase2b
mpirun -np 1 -x ELPA_FORCE_gpu=1 thermal_anharmonicity_fase2b
\end{verbatim}
The kernel names are those printed by \file{elpa2\_print\_kernels}. For
the GPU use \texttt{ELPA\_FORCE\_gpu}: the option \texttt{nvidia-gpu} has a
dash, which an environment variable cannot have. ELPA does not always
confirm a forced option in the output; for the GPU, \texttt{nvidia-smi}
shows whether the process uses it.

\paragraph{Suggestions.}
\begin{itemize}
\item Keep \opt{elpablock} at 64 or below. ELPA's two-stage solver gave
  wrong eigenvalues, without any error, with blocks of 300 on polyads of
  3400 states and more; ELPA's own test on the same matrix failed in the
  same way. The GPU kernel needs a power of 2.
\item The kernel of the two-stage solver only enters the
  back-transformation of the eigenvectors: changing it changes the speed,
  not the eigenvalues.
\item For polyads of up to a few thousand states one MPI process is the
  fastest: with more processes, distributing the matrices costs more than
  it gains (Chapter~\ref{ch:tut-pyrene}). Run several starting states as
  independent runs instead.
\item The GPU helped a little with one process (pyrene tutorial on a
  laptop with a GeForce GTX 1650 Ti: 1:50 instead of 2:20), and not at all
  with four. Consumer GPUs are slow in double precision. Several
  processes sharing one GPU have not been tested yet.
\item After installing a new ELPA, or on a new machine, compare one run
  with \opt{eigensolver}~=~2 (\texttt{pdsyevx}): the spectra must agree
  to rounding.
\end{itemize}

\section{Checking the build}

The directory \file{tests/checks} contains checks that compare the
code with exact reference implementations; each \file{run.sh} ends with
\texttt{OK}. They were designed and written with the support of Claude,
the AI model by Anthropic, during the 2026 review of the code. Phase 1
alone can also be built without MPI, PETSc or PARI, with
\file{tests/checks/fase1\_standalone/build.sh}. This is useful
on a laptop, to inspect the resonances before running phases 2a and 2b on
a cluster.
